% GENERATED, edit the French unit or the translation table
\mxunit{seconde-s05}{Term review and remediation}\mxsubtitle{16--20 November 2026}
\section{Individual review}
\begin{method}[Instructions]
Work without the course notes for 25 minutes, then check with the methods in the booklet. For each question, note: solved alone, solved with help, or to redo. Justify the answers.
\end{method}
\begin{example}[Review of sequences 1 to 4]\begin{enumerate}
\item Reduce $90/126$ to lowest terms. Prove that the sum of two odd integers is even.
\item With $A(-1;2)$ and $B(3;-2)$, calculate $\vect{AB}$, $AB$ and the midpoint of $[AB]$ in an orthonormal frame.
\item Determine an equation of $(AB)$; does $C(5;-4)$ belong to this line?
\item Calculate $[-2;4]\cap]1;7]$ and $[-2;4]\cup]1;7]$.
\item Solve $|x-1|\leq4$ and explain why $1/3\notin\D$.
\end{enumerate}\end{example}
\mxanswerspace[9]{1. $\pgcd(90,126)=18$, so $90/126=5/7$. If $u=2k+1$ and $v=2\ell+1$, then $u+v=2(k+\ell+1)$ is even.\\2. $\vect{AB}=(4;-4)$; $AB=4\sqrt2$; midpoint $I(1;0)$.\\3. $y=-x+1$. For $C$, $-4=-5+1$: it belongs to $(AB)$.\\4. Intersection $]1;4]$; union $[-2;7]$.\\5. $x\in[-3;5]$. Writing $1/3=a/10^n$ would require $10^n$ to be divisible by $3$, which is impossible.}
\begin{teachernote}Possible marking out of 20: 4 points per question. Assess calculation skill, choice of method and quality of justification separately.\end{teachernote}
\mxsection{Correcting an error and revisiting a method}
\begin{example}[Analysis of student work]Correct: (a) ``$\vect{AB}=(x_A-x_B;y_A-y_B)$''; (b) ``The zero vector directs every line''; (c) ``$|x-2|\leq3$ gives $x\in[2;3]$''; (d) ``$1/3\approx0.333$, so $1/3$ is decimal''.\end{example}
\mxanswerspace[5]{(a) Subtract the origin from the endpoint: $(x_B-x_A;y_B-y_A)$. (b) A direction vector must be nonzero. (c) The distance to $2$ is at most $3$: $x\in[-1;5]$. (d) $0.333$ differs from $1/3$: $3\times0.333=0.999\ne1$. A decimal approximation does not determine the nature of the approximated number.}
\begin{example}[Remediation track]\begin{enumerate}
\item Reduce $72/108$, then $\frac34-\frac16$.
\item With $P(1;1)$ and $Q(5;3)$, determine the midpoint, a direction vector and an equation of $(PQ)$.
\item Determine $]-\infty;3]\cap[-1;+\infty[$, then solve $|x+2|<1$.
\end{enumerate}\end{example}
\mxanswerspace[5]{1. $2/3$; $7/12$. 2. Midpoint $(3;2)$, direction vector $(2;1)$, equation $x-2y+1=0$. 3. Intersection $[-1;3]$; solution of the inequality $]-3;-1[$.}
\mxsection{Applying and self-assessing}
\begin{example}[Synthesis problem]Two displacements are represented by $\vec u(6;4)$ and $\vec v(9;6)$.
\begin{enumerate}\item Show that they are collinear and find the coefficient relating $\vec v$ to $\vec u$.
\item The starting point is $A(-2;1)$. Determine the endpoints $B$ and $C$ associated with the displacements $\vec u$ and $\vec v$.
\item Find an equation of the line through these three points.
\item A point $M(x;y)$ of this line must satisfy $|x-4|\leq3$. Give the interval of possible abscissae.
\end{enumerate}\end{example}
\mxanswerspace[7]{1. $\vec v=\frac32\vec u$. 2. $B(4;5)$ and $C(7;7)$. 3. A direction vector is $(3;2)$; the equation is $2x-3y+7=0$, satisfied by the three points. 4. $|x-4|\leq3\iff x\in[1;7]$.}
\begin{keypoint}[My review]
My priority: \dotfill\par The mistake to avoid: \dotfill\par The method to reuse: \dotfill\par An exercise to redo without help: \dotfill
\end{keypoint}
\begin{teachernote}Validate the remediation with a similar exercise done without help after a delay. Orally, ask for a sentence linking the calculation to the geometric conclusion.\end{teachernote}
